% =============================================================================
% Front matter: Provisional Abstract
% Input-ready fragment for the lean toy-model research proposal.
% This file intentionally contains no preamble or document environment.
% =============================================================================

\chapter*{Provisional Abstract}
\addcontentsline{toc}{chapter}{Provisional Abstract}
\label{ch:provisional-abstract}

This proposal defines a falsifiable research program for a classical
one-medium brane--bulk toy analog. The model assumes preferred time and four
Euclidean spatial dimensions. Observed three-dimensional space is identified
with an ordered, finite-thickness, shear-supporting brane within a larger
compressible medium; the brane and bulk are ordered and de-structured states of
the same globally conserved substance rather than separate materials. Light is
sought as a guided transverse brane excitation with two polarizations, while
particles are hypothesized to be finite oriented throats supported in part by a
trapped transverse standing mode. Throat orientation supplies the proposed
electric sign. Localized one-way throat drainage, throat-region pressure,
stress, and order conversion, the separate distributed reordering of bulk
material into ordered brane material represented by the positive bulk-to-brane
contribution to \(S_{\mathrm{leak}}\), and the resulting reservoir response are
proposed as the material basis for a calibrated gravity observable. Motion and
acceleration of the same complete dressed configuration are intended to produce
magnetism, radiation, passive response, and inertia.

The central question is whether one mathematically coherent medium can support
all required sectors simultaneously without independent retuning. The research
therefore begins with a far-field-first capability analysis rather than with a
premature nonlinear particle solve. Guided light and the stronger localized
photon-packet hypothesis are tested alongside gravity, static electricity,
magnetism, radiation from accelerating charge, inertia, drag, leakage, and unwanted
modes. Each calculation must determine what behavior a brane observer would
see, what mechanism and response operator could produce it, and what fields,
coefficients, characteristic speeds, source conventions, conservation laws,
and exclusions it imposes on the shared medium. The resulting requirements are
then intersected rather than judged as separately tuned successes.

If that intersection is nonempty, the surviving field content, dynamical
branch, constitutive terms, parameters, boundary data, constraints, and
reservoir variables will be frozen as Candidate Parent System V1. The candidate
must support a stable finite brane and complete background spectrum, after
which all far-field sectors are rederived from that same system and parameter
region. Only a candidate that survives this common-medium test advances to the
nonlinear program: self-consistent supported throats, orientation partners and
species structure, moving-throat response, two-body forces, finite-thickness
contact and possible annihilation or neutral binding, and the fixed point
between localized one-way throat drainage, the global bulk reservoir, and the
separate distributed \(S_{\mathrm{leak}}\) return across the porous brane.

Success would establish a controlled classical analog architecture within its
stated scope, not a derivation of the real universe, exact general relativity,
exact electromagnetism, or quantum theory. An empty compatibility region, an
unavoidable extra mode or loss channel, failure of the supported-throat branch,
or inability to close the global conservation and reservoir ledgers would be an
equally legitimate result because it would identify where the proposed
architecture fails.
